\section{Validation}
\label{sec:valid}

\begin{guidance}[Validation and automated test harness (15 marks)]
Show \textbf{passing tests for every required case} and \textbf{screenshots or a
short video} of the working system. The brief's required tests (p.~4):
CONNECT returns OK and NEW returns an empty board; a legal move changes the
board exactly as expected; an occupied-cell move returns NACK:OCCUPIED and
leaves the board unchanged; an out-of-range move returns NACK:RANGE; a move
after game over returns NACK:NOGAME (or your documented equivalent); scripted
win and draw sequences return the correct result; and, if you built an AI
opponent, it never makes an illegal move. Map each one to a named test.
\textbf{Report results on the real board}, and say clearly which results come
from a simulator, if you used one.
\end{guidance}

\subsection{One fixture puts the board into a known state before every test}
The harness uses pytest~\cite{pytest} and pySerial~\cite{pyserial}. The port
name comes from the \texttt{SERIAL\_PORT} environment variable, and the suite
skips if it is unset, so it cannot fail on a machine without a board. Before each
test the fixture sends \texttt{CONNECT} until the board answers \texttt{OK}
(for up to 10~s, so it does not depend on the boot banner or on a reset when the
port opens), waits for stale replies, and sends \texttt{DISCONNECT}. Every test
therefore starts in IDLE with no session. Run the core suite with
\texttt{python -m pytest -v test\_protocol.py}, the extension suite with
\texttt{test\_extras.py}, and the whole suite with a checksum on every request by
setting \texttt{USE\_CRC=1}. The board results in this section come from the real
board on \texttt{COM14} (Windows, Python 3.11.9, pytest 7.4.3). Results from a PC
build of the firmware are labelled as such.

\subsection{21 of 23 tests passed on the first board run, and the two failures were real}
The core suite (\texttt{test\_protocol.py}) has 23 tests: the required cases of Table~\ref{tab:required}, the state
machine (\texttt{NOCONN}, \texttt{NOGAME}, \texttt{NOTURN}, resign, disconnect)
and the robustness cases of Table~\ref{tab:threats}. On the first run on the
board, 21 passed and 2 failed in 12.45~s; the extension tests did not exist yet. In \texttt{test\_out\_of\_range\_rejected}
the first four indices (9, 10, 99, $-1$) returned \texttt{NACK:RANGE}, but the
fifth, a 15-digit number, returned \texttt{NACK:BADCMD}. In the oversize test a
103-character line returned \texttt{NACK:BADCMD} instead of
\texttt{NACK:TOOLONG}. Both failures involve frames much longer than the
passing 6-character commands, which led us to the receive path described in
Section~\ref{sec:impl}. The second run, after we replaced the HAL receive path
with the register-level handler, is \ph{result: N passed, M failed in T s}.

\begin{table}[t]
\centering
\caption{The brief's required tests and their results on the board.}
\label{tab:required}
\scriptsize
\begin{tabular}{@{}L{3.2cm}L{3.0cm}cc@{}}
\toprule
Requirement (brief, p.~4) & Test name & Run 1 & Run 2 \\
\midrule
\texttt{CONNECT} returns \texttt{OK}; \texttt{NEW} returns an empty board & \texttt{test\_connect\_returns\_ok}; \texttt{test\_new\_returns\_empty\_board} & pass & \ph{res} \\
A legal move changes the board exactly as expected & \texttt{test\_legal\_move\_changes\_board\_exactly} & pass & \ph{res} \\
Occupied cell returns \texttt{NACK:OCCUPIED}, board unchanged & \texttt{test\_occupied\_move\_rejected\_and\_board\_unchanged} & pass & \ph{res} \\
Out-of-range move returns \texttt{NACK:RANGE} & \texttt{test\_out\_of\_range\_rejected} & fail & \ph{res} \\
Move after game over returns \texttt{NACK:NOGAME} & \texttt{test\_move\_after\_game\_over\_rejected} & pass & \ph{res} \\
Scripted win and draw give the correct result & \texttt{test\_scripted\_win\_x}; \texttt{test\_scripted\_win\_o}; \texttt{test\_scripted\_draw} & pass & \ph{res} \\
AI never makes an illegal move & \texttt{test\_ai\_never\_makes\_an\_illegal\_move\_and\_never\_loses} (\texttt{test\_extras.py}) & n/a & \ph{res} \\
\bottomrule
\end{tabular}
\end{table}

Table~\ref{tab:required} shows that the required behaviour was correct on the
first run except for very long indices. The out-of-range test failed only on its
last value; \texttt{NACK:RANGE} for 9, 10, 99 and $-1$ was already correct.

\subsection{24 extension tests check the checksum, the AI and undo}
Table~\ref{tab:ext} summarises \texttt{test\_extras.py}. The AI test plays 20 seeded
random games and checks that every reply differs from the previous board by
exactly the human's X and one new O, and that the human never wins. The checksum
tests include the substitution of \texttt{MOVE:4} by \texttt{MOVE:5} with the old
checksum, and the known CRC-8 value. Setting \texttt{USE\_CRC=1} reruns all 47
tests with a checksum on every request and verifies every reply.

\begin{table}[t]
\centering
\caption{Extension tests and their results on the board.}
\label{tab:ext}
\scriptsize
\begin{tabular}{@{}L{1.45cm}L{0.65cm}L{4.15cm}L{1.1cm}@{}}
\toprule
Group & Tests & What they check & Board \\
\midrule
Checksum & 9 & test vector, reply carries a suffix only if the request did, wrong checksum rejected, corrupted payload caught & \ph{res} \\
AI & 7 & one AI move per human move, blocks a threat, takes a win, 20 random games with no illegal AI move and no human win & \ph{res} \\
History, undo & 8 & move list, undo restores board and turn (both modes), \texttt{NOUNDO}, \texttt{NOGAME}, \texttt{NOCONN} & \ph{res} \\
All 47, \texttt{USE\_CRC=1} & 47 & the whole suite with a checksum on every request & \ph{res} \\
\bottomrule
\end{tabular}
\end{table}

\ph{Optional, keep only if you run the PC build (tools/hostsim) yourself: before flashing
we also ran all 47 tests against a PC build of \texttt{main.c} with the HAL stubbed out,
plain and with \texttt{USE\_CRC=1}, and all passed. This is not a board result.}

\begin{figure}[t]
\centering
\figph{5cm}{Screenshot of your GUI connected to the board, e.g. mid-game or
just after a rejected move, with the message the user sees.}
\caption{\ph{What the screenshot shows and what to notice in it.}}
\label{fig:screenshot}
\end{figure}

\ph{Refer to Fig.~\ref{fig:screenshot} here: one sentence on what the reader
should notice, e.g.\ the dialog for a rejected move.}

\subsection{Code review found five bugs before the board did, and the board found one more}
Reading the first firmware against the brief found five defects: the USART1
interrupt was never enabled, so no command was received; \texttt{atoi} turned
\texttt{MOVE:abc} into a move on cell 0; the tail of an oversized line ran as a
command; \texttt{NEW} was accepted in the middle of a game; and the hand-written
boards in the first draw test were wrong, so that test would have failed on
correct firmware. We fixed these before the first board run.
The board then exposed the receive-path overrun described in
Section~\ref{sec:impl}. \ph{Add any bug you hit while running the board or the GUI
that is not listed here.}
